\scanpage{20}
SYSTÈMES de RACINES (15) \hfill Tribu n° 59

\uline{n° 5} - Identifier $V'$ à son dual à la p. 195. A la p. 196 décrire plus explicitement $P(R)$ et $Q(R)$ : $Q(R)=L_0\cap V'$, $P(R)$ est la proj. orthogonale de $L_0$ sur $V'$, (cf. résultat ajouté au § 6, n° 8) ; on a $P=Q+Zv'$ où $v'$ est la projection de $(1,0,\ldots,0)$, (on a $(\ell+1)v'=(\ell,-1,\ldots,-1)$.

En bas de la p. 196, ne pas identifier tout de suite $W$ et le groupe d'autom. de $R^{\ell+1}$ d'où il peut venir, mais remonter $W$ en un groupe $G$ ; déterminer $S(R^{\ell+1})$, $S(R^{\ell+1})^G$, $S(V')$ et $S(V')^W$. Diviser par la factorielle à la l.15, p. 197.

Dans (XI) dire, tout de go, que l'autom. échange $i$ et $\ell-i+1$. Ajouter l'effet de $\operatorname{Aut}(R)$ sur $P/Q$.

\uline{n° 6} - Prendre $\ell\geq1$... Pas de dual à la p. 199. Dire seulement que $\gamma=(4\ell-2)(\ell+1)$, en référant à la formule (13)...
Dans (IX) et (X) dire que $W(R)$ se compose des applic. linéaires telles que... ; c'est seulement un produit semi-direct. Diviser par la factorielle dans (7), p. 201. Remplacer ``Etant'' par ``Comme ils sont'', (p. 201, l.6). Ce qu'on dit dans (XI) concerne le cas $\ell\geq2$ ; ajouter que la même conclusion est évidente pour $\ell=1$.

\uline{n° 7} - Prendre $\ell\geq1$. Parler de ``base'' d'un système de racines. Dans (V) mettre un 4 au dénominateur de la forme canonique. Répéter les deux réseaux dans (VIII).

\margin{--}

\uline{n° 8} - Prendre $\unsure{l|1}\geq2$, et dire que c'est réductible pour 2. Dire que $R^+$ est l'ensemble des $\varepsilon_i\pm\varepsilon_j$ pour $i>j$, (en effet on a $R=R^+\cup(-R^+)$). Dans (V) on a $R=R$.

\mspage{19}
